\part{Foundations}

\chapter{An Introduction to Mathematical Reasoning}
\chapter{Sets}

Almost all of modern mathematics can be put in the context of the study of \tib{sets} and their \tib{elements}. Generally, given a collection of objects $a,b,c,d,\hdots$, they can be taken together as a set, generally written in curly brackets and represented by a capital letter:
\begin{align*}
	M := \lr{a,b,c,d,e,\hdots}
\end{align*}
However, an approach where any collection of mathematical objects forms a set quickly runs into trouble - famously, \tib{Russel's paradox} asks:
\begin{quote}
	\tit{"Let $R$ be the set of all sets that don't contain themselves. Does $R$ contain itself?"}
\end{quote}
If $R$ doesn't contain itself, then by definition it must contain itself, but if it contains itself, then it cannot contain itself! If $R \in R$ is true, then it must be false! We arrive at a classic paradox. The conclusion, then, is that the set of all sets that don't contain themselves cannot exist - we cannot define sets to be any collection of things defined by any property, and must instead establish ground rules about which sets exist and exactly how we can use sets we already know to construct new ones.
\newpar
These ground rules are given in the form of \tib{axioms}, which are statements that mathematicians take for granted, which can then be used as a starting point to then derive "the rest of mathematics" from. The following chapters contain a brief and hopefully easily understandable introduction to the ZFC axiom system, which was developed in the early 20th century by Ernst Zermelo and Abraham Fraenkel and is now the most widely accepted axiomatization of mathematics.
\section{Infinity: $\emptyset$ and $\bN$}
\section{Extensionality}
To say "$b$ is an element of $M$", or equivalently "$b$ contains $M$", we simply write:
\begin{align*}
	b \in M
\end{align*}
One of the most basic axioms of set theory is the \tib{Axiom of Extensionality}, which states:
\begin{axiom}{Axiom of Extensionality}{}
	Two sets $A$ and $B$ are equal if and only if they contain the same elements.
\end{axiom}
The Axiom of Extensionality establishes that sets are uniquely defined by the objects they contain and have no additional structure. In particular:
\begin{enumerate}
	\item The order of elements in a set doesn't matter. $\lr{1,2}$ and $\lr{2,1}$ are the same set.
	\item The number of times an element appears within the same set doesn't matter. $\lr{1}, \lr{1,1}$, and $\lr{1,1,1}$ are the same set.
\end{enumerate}
\begin{importantdefinition}{Subsets}{}
	We call a set $A$ a \tib{subset} of a set $B$ if and only if every element of $A$ is contained in $B$. We write "$A$ is a subset of $B$" as $A \subseteq B$.
\end{importantdefinition}
The demand that "every element of $A$ is contained in $B$" can be stated more formally as "$x \in A$ implies $x \in B$". Importantly, this means that every set is a subset of itself - we always have $A \subseteq A$.

\newpar
This definition gives us another way of phrasing the Axiom of Extensionality - two sets are equal if and only if each is a subset of the other:
\begin{align*}
	A = B \Leftrightarrow A \subset B \tn{ and } B \subset A
\end{align*}
This idea is incredibly important for writing out proofs in practice - if a problem asks you to establish that two sets $A$ and $B$ are the same, then the easiest approach is almost always to show that any element of $A$ must be contained in $B$, and to then show that any arbitrary element of $B$ must be contained in $A$.
\newpar
The subset relation is frequently written as $A \subset B$. However, other authors reserve $A \subset B$ to mean "$A$ is a subset of $B$ that isn't equal to $B$" - we call such a set a \tib{proper subset of $B$}. I will therefore attempt to only use the symbols $\subseteq$ and $\subsetneq$, which clear up this ambiguity.
\newpar
\section{Specification}
\section{Pairing, Union, and Power Sets}
\section{Relations}
\section{Functions}
\section{Cartesian Products}
\begin{importantdefinition}{General Products in Category Theory}{}
	Let $C$ be a category. Let $X_1$ and $X_2$ be objects of $C$. Then the \tbf{product of $X_1$ and $X_2$} is an object $X_1 \times X_2$ equipped with two morphisms $\pi_1 : X \to X_1$, $\pi_2 : X \to X_2$ with the property that for every object $Y$ and every pair of morphisms $f_1 : Y \to X_1$, $f_2 : Y \to X_2$, there exists a unique morphism $f : Y \to X_1 \times X_2$ such that:
	\begin{enumerate}
		\item $f_1 = \pi_1 \circ f$
		\item $f_2 = \pi_2 \circ f$
	\end{enumerate}
	We call this property the \tbf{universal property of product objects.}
\end{importantdefinition}
In the case of sets, this definition leads us to the familiar cartesian products, where $\pi_1$ and $\pi_2$ are the projections to the first and second component of the product respectively. We can generalize the cartesian product to uncountably infinite products as follows:
\begin{importantdefinition}{Cartesian Product of sets}{}
	Let $I$ be any set. Let $(X_i)_{i \in I}$ be a Family of nonempty sets indexed by $I$, i.e. there exists a function assigning to each element $i \in I$ an element of $(X_i)_{i \in I}$. Then the cartesian product of $(X_i)_{i \in I}$ is the unique set:
	\begin{align*}
		\prod_{i \in I} X_i := \lr{x : I \to \bigcup_{i \in I} X_i \mid x_i := x(i) \in X_i \tn{ for all } i \in I}
	\end{align*}
\end{importantdefinition}
Written more simply: Each element of $\prod_{i \in I} X_i$ is a function assigning to each $i \in I$ an element $x(i) \in X_i$, where we generally write the argument in the index ($x_i := x(i)$), to stay consistent with the established notation in the countable case. 
\newpar
This definition is simply a reinterpretation of the definition of finite cartesian products: It amounts to viewing, for example, the tuple $ T := (x,y,z)$ as a map:
\begin{align*}
	T : \lr{1,2,3} &\to \lr{x,y,z}
\end{align*}
such that $T(1) = x$, $T(2) = y$, and $T(3) = z$. 
\begin{importantdefinition}{Projections}{}
	Let $J \subset I$. Then we define the projection of $I$ to $J$ as the map
	\begin{align*}
		\pi_J := \pi_{J}^I : \prod_{i \in I} X_i &\to \prod_{j \in J} X_j\\
		x &\mapsto \lr(x|_J : J \to \bigcup_{j \in J} X_j)
	\end{align*}
	which restricts the domain of a given element $x$ of the product to the subset $J$.
\end{importantdefinition}
In particular, if $J$ is a set consisting of a single element $j$, we have:
\begin{align*}
	\pi_j := \pi_{\lr{j}}^I : \prod_{i \in I} X_i &\to X_j\\
	x &\mapsto x_j
\end{align*}
\section{Replacement}
\section{Regularity and Well-Foundedness}
\section{Choice}
\section{Cardinality}
